\section{Excluding six or more pieces}\label{sec:outline-classes}

By \cref{prop:partition}, \cref{thm:at-most-five} follows from one theorem for each of the five
classes. The five exclusions use different mechanisms. The classes $\cls{R}$ and $\cls{M3}$ are
excluded for every normal vertex period, with no reference to the cost: their obstructions are
local and come from the maximum principle alone. The classes $\cls{M1}$, $\cls{F}$ and $\cls{M2}$
are excluded only for competitive periods, through the Casimir cap and estimates over the whole
period. This section states the five theorems and explains their proofs, which are in
\cref{sec:reverse,sec:short-runs,sec:forward,sec:mixed,sec:smaller-short}.

\subsection{Classes R and M3}\label{sec:outline-reverse}

The two classes whose obstruction is local are an entirely reverse period and a mixed period with
a maximal reverse run of length other than two. Both exclusions rest on two facts about a few
consecutive pieces, which hold for every normal vertex period, whatever its cost and its number of
pieces. The first says that across three reverse switches the active margin of an interior piece
cannot be a minimum of the four.

\begin{lemma}[Reverse-four principle]\label{lem:reverse-four}
Let four consecutive pieces $k,\dots,k+3$ of a normal vertex period have three reverse switches
between them. Then
\begin{enumerate}[label=\textup{(\alph*)}]
\item it is not the case that $a_{k+2}\le a_k$, $a_{k+2}\le a_{k+1}$ and $a_{k+2}\le a_{k+3}$;
\item it is not the case that $a_{k+1}\le a_k$, $a_{k+1}\le a_{k+2}$ and $a_{k+1}\le a_{k+3}$.
\end{enumerate}
\end{lemma}

It is proved from two exact linear identities between the switching functions at the three
switches (\cref{lem:rev-identities}) and the sign of an explicit cost derivative
(\cref{lem:rev-N}). The second fact says that across two switches that return to the same label,
the active margins are strictly log-concave; it is read off the explicit switching function on the
middle piece.

\begin{lemma}[Three-arc log-concavity and turn ordering]\label{lem:three-arc}
Let $k,k+1,k+2$ be consecutive pieces of a normal vertex period.
\begin{enumerate}[label=\textup{(\alph*)}]
\item If $i_{k+2}=i_k$, then $a_ka_{k+2}<a_{k+1}^2$; in particular $a_{k+1}$ is not at most both
  $a_k$ and $a_{k+2}$.
\item If the switch $k$ is reverse and the switch $k+1$ forward, then $a_k<a_{k+2}$; if the switch
  $k$ is forward and the switch $k+1$ reverse, then $a_{k+2}<a_k$.
\end{enumerate}
\end{lemma}

An entirely reverse period has a global minimum of its active margins, and the four pieces around
it have three reverse switches between them; this contradicts \cref{lem:reverse-four}.

\begin{theorem}[Class $\cls{R}$]\label{thm:class-R}
No normal vertex period has every switch reverse.
\end{theorem}

In a reverse run of a mixed period, the two lemmas leave room for exactly one configuration: a run
of two switches around a strict minimum.

\begin{proposition}[Run structure]\label{prop:run-structure}
Let $P$ be a mixed normal vertex period, and let the switches $k,\dots,k+L-1$ form a maximal
reverse run. Then $L=2$, and the active margin of the centre piece $k+1$ is strictly smaller
than those of the four pieces around it:
\[
  a_{k+1}<a_{k-1},\qquad a_{k+1}<a_k,\qquad a_{k+1}<a_{k+2},\qquad a_{k+1}<a_{k+3}.
\]
\end{proposition}

\begin{theorem}[Class $\cls{M3}$]\label{thm:class-M3}
No normal vertex period is mixed with a maximal reverse run of length other than two.
\end{theorem}

\Cref{thm:class-M3} is the first conclusion of \cref{prop:run-structure}. The second conclusion,
the strict minimum at every centre, is the input of the remaining mixed classes, and of the proof
of ellipticity (\cref{prop:ellipticity}). The proofs are in \cref{sec:reverse}, and the identities
behind them in \cref{app:reverse-estimates}.

\subsection{Class M1}\label{sec:outline-short-runs}

A period of class $\cls{M1}$ is mixed, every maximal run of reverse switches has length two and
every maximal run of forward switches has length one. Its cyclic sign word is therefore
$(--+)^m$: every reverse run is a pair of reverse switches around a centre piece, and consecutive
pairs are separated by single forward switches.

\begin{theorem}\label{thm:class-M1}
  No competitive normal vertex period belongs to class \cls{M1}.
\end{theorem}

The theorem uses neither the compact domain nor a lower bound on the number of pieces. Its proof
compares two ways of counting how far the switch vectors $y_k$ turn in one period, measured in a
positive definite form $S$ built from the normal multiplier. At a turn the switch vector only
changes sign, so all the turning happens at the $m$ centres. A whole-period winding identity fixes
the total turning at exactly $(m+1)\pi/3$, because the frame of a competitive period turns by
exactly $\pi/3$ (\cref{prop:m1-winding}). A local estimate at each centre, which combines the
stationarity identity, the Gram identity and the Casimir cap, shows that each centre turns the
switch vectors by more than $2\pi/5$ (\cref{prop:m1-angle-gap}). The average $(m+1)\pi/(3m)$ is at
most $2\pi/5$ as soon as $m\ge5$, and the sign sum (\cref{lem:sign-sum}) leaves $m=2$ as the only
other case. With two centres the recurrence of the switch vectors closes up only in a rigid
configuration whose cost can be computed exactly, and it exceeds $\Aoct$
(\cref{prop:m1-two-centre}). The proof is in \cref{sec:short-runs}.

\subsection{Class F}\label{sec:outline-forward}

The smoothed octagon's period is entirely forward. No other entirely forward period with six or
more pieces is competitive.

\begin{theorem}[Class F]\label{thm:class-F}
No competitive normal vertex period with at least six pieces has every switch forward.
\end{theorem}

The proof plays two consequences of competitiveness against each other. On one side, the spectral
cost--Casimir inequality gives the sharper cap
$\gamma^2\le(3\Aoct/\pi-\frac32)^2/12$
(\cref{prop:fwd-spectral,cor:fwd-cap}). On the other side, $\gamma^2$ can be
read off at any piece from the two switch norms on either side of it and their cross pairing
(\cref{lem:gram}), and we read it at a piece of smallest active margin. In an entirely forward
period $N\equiv1\pmod3$, so a period with $N\ge6$ has at least seven pieces, and its cost is a sum
of seven or more edge terms that depend only on consecutive active margins. At the octagon's
margins each edge costs exactly a quarter of $\Aoct$, so seven edges fit into the budget only if
the smallest margin is markedly smaller than the octagon's. The switching conditions at the five
pieces around that small margin then hold the two switch norms close to their least possible value
$3/32$, and with such norms the Gram identity makes $\gamma^2$ slightly larger than
competitiveness allows. The central pairing retains a positive logarithmic correction; the resulting
vertex bounds for the Casimir exceed the spectral cap (\cref{fig:forward}). The proof is in \cref{sec:forward}, and its scalar estimates in
\cref{app:forward-estimates}.

\subsection{Class M2}\label{sec:outline-mixed}

A period of class $\cls{M2}$ has switches of both signs, every maximal reverse run has exactly two
switches, and some maximal forward run has at least two. Unlike the other classes, such a period
has no rigid sign word: the forward runs may have any lengths, the number of reverse runs is
arbitrary, and the factors are unrelated.

\begin{theorem}[Class $\cls{M2}$]\label{thm:class-M2}
Let $P$ be a competitive normal vertex period whose state stays in $\hstarstar$ and which has at
least six pieces. Then the sign sequence of $P$ is not of class $\cls{M2}$.
\end{theorem}

The exclusion is therefore local. We look at a window of at most seven consecutive pieces around
the smallest active margin of the period, read the normal multiplier of the whole period in that
window, and derive a finite list of scalar necessary conditions, which are then shown to be
inconsistent. The whole period enters only through three inputs: the compact domain, the Casimir
cap, and a lower bound on the Casimir obtained from the winding of the switch vectors over the
period. The first step locates the window.

\begin{proposition}[The minimum is a reverse centre]\label{prop:m2-centre}
Let $P$ be a competitive normal vertex period of class $\cls{M2}$ with at least six pieces. There is
an index $m$ with $a_m\le a_k$ for every $k$ such that the switches $m-1$ and $m$ are reverse and the
switches $m-2$ and $m+1$ are forward.
\end{proposition}

Of the two turns beside that centre we choose the one with the smaller active margin; after
reflecting the period and reversing time if necessary, it follows the centre
(\cref{lem:m2-reflection}). The forward run that starts at this turn has one switch or more.

If it has at least two, the run is \emph{nonshort}, and the switch after the next piece is forward
or reverse (\cref{prop:m2-window,prop:m2-continuation}). A comparison that moves one margin to a
boundary value turns the window into one scalar residual $N$ that must be negative
(\cref{lem:m2-comparison}), together with a short list of necessary bounds
(\cref{lem:m2-necessary}). A reverse successor contradicts $N<0$ directly (\cref{thm:m2-reverse}).
A forward successor leaves a small region of parameters on which the Casimir exceeds the cap
(\cref{thm:m2-forward}). \Cref{sec:mixed} gives this part of the proof, and
\cref{app:nonshort-estimates} its polynomial inequalities.

If the forward run has a single switch, it is a short bridge between two reverse runs. Its data
are fixed by six real numbers, and no free continuation of the period is left to compare with. The
proof works with a fixed list of necessary conditions on those numbers, the \emph{smaller-short
window} (\cref{def:ss-window}), and shows that the list is inconsistent
(\cref{thm:m2-smaller-short}). It proceeds through seven steps, which we call gates. Gates S0 to
S3 confine the shape of the bridge to a small box, gate S4 gives a lower bound for the Casimir,
gate S5 shows that the bridge is sheared in a definite direction, and gate S6 shows that a bridge
sheared in that direction costs more, together with the rest of the period, than the octagon.
\Cref{sec:smaller-short} gives this part of the proof, and \cref{app:short-estimates} its scalar
estimates.

\subsection{Proof of \texorpdfstring{\cref{thm:at-most-five}}{the bound on the number of pieces}}\label{sec:outline-five-proof}

\begin{proof}[Proof of \cref{thm:at-most-five}]
Suppose $N\ge6$. By \cref{prop:partition} the period is in one of the five classes, and each
class is excluded by a theorem proved in the following sections:
\begin{center}
\begin{tabular}{@{}lll@{}}
\toprule
Class & Excluded by & Hypotheses used\\
\midrule
$\cls{R}$ & \cref{thm:class-R} & none\\
$\cls{M3}$ & \cref{thm:class-M3} & none\\
$\cls{M1}$ & \cref{thm:class-M1} & competitive\\
$\cls{F}$ & \cref{thm:class-F} & competitive, $N\ge6$\\
$\cls{M2}$ & \cref{thm:class-M2} & competitive, state in $\hstarstar$, $N\ge6$\\
\bottomrule
\end{tabular}
\end{center}
Only the classes $\cls{F}$ and $\cls{M2}$ use the bound on the number of pieces.
\end{proof}

The classes $\cls{R}$ and $\cls{M3}$ are excluded for every normal vertex period, with no
reference to the cost: their obstructions are local in the period and come from the maximum
principle alone. The classes $\cls{M1}$, $\cls{F}$ and $\cls{M2}$ are excluded only for
competitive periods, through the Casimir cap and estimates over the whole period.
\Cref{sec:outline-terminal} treats the periods with at most five pieces.
